
\section{Politica y Compromiso de la Dirección}
SKYFILMS reconoce que la seguridad operacional es un valor fundamental y prioritario en todas nuestras actividades. La organización se compromete a demostrar un liderazgo activo y visible en la promoción de una cultura de seguridad positiva en toda la organización. Este compromiso se manifiesta a través de los siguientes principios que serán revisados anualmente:

\begin{enumerate}
    \item \textbf{Mejora Continua:} SKYFILMS se compromete a mejorar continuamente el nivel de rendimiento en materia de seguridad operacional, buscando activamente oportunidades para fortalecer nuestros procesos y reducir los riesgos.
    \item \textbf{Cultura de Seguridad:} Fomentar y mantener una cultura de seguridad operacional positiva en SKYFILMS, donde cada empleado se sienta responsable y partícipe en la identificación y gestión de riesgos.
    \item \textbf{Cumplimiento Normativo:} Cumplir con todos los requisitos normativos aplicables, incluyendo el RAC 100 y 219, así como cualquier otra regulación que afecte nuestras operaciones.
    \item \textbf{Recursos Adecuados:} Proporcionar los recursos humanos, financieros y técnicos necesarios para garantizar la seguridad de nuestras operaciones y el correcto funcionamiento del SMS.
    \item \textbf{Responsabilidad de la Gestión:} Asegurar que la seguridad operacional sea una responsabilidad primordial de todos los administradores y supervisores en SKYFILMS.
    \item \textbf{Comprensión e Implementación:} Garantizar que esta política de seguridad operacional sea comprendida, implementada y mantenida en todos los niveles de la organización.
    \item \textbf{Comportamiento inaceptable:} SKYFILMS define claramente los tipos de comportamientos que son inaceptables en relación con nuestras operaciones. Estos incluyen, pero no se limitan a:
    
        \begin{enumerate}[label=(\roman*)]
            \item Violaciones deliberadas de las regulaciones, políticas y procedimientos de seguridad.
            \item Conductas negligentes que pongan en riesgo la seguridad de las operaciones.
            \item Falsificación de registros o informes relacionados con la seguridad.
            \item Interferencia con los procedimientos de control de seguridad.
            \item Operar bajo la influencia de alcohol o sustancias psicoactivas.
        \end{enumerate}
    
    SKYFILMS se reserva el derecho de aplicar medidas disciplinarias proporcionales a la gravedad o falla en caso de comportamientos inaceptables. Sin embargo, no se aplicarán medidas en aquellos casos en que un empleado realice un informe de seguridad de buena fe, incluso si dicho informe revela una infracción no intencional de los procedimientos o regulaciones, siempre y cuando la acción no constituya negligencia grave o una violación deliberada.

\end{enumerate}

\vspace{1cm}

\begin{center}

\includegraphics[width=4cm]{Imagenes/firma.jpg}

\vspace{-1.2cm}

\rule{7cm}{0.5pt}

\vspace{0.3cm}
\textbf{Andrés Felipe Pérez Sierra}\\
C.C. 1.108.206.142\\
Gerente General\\


\end{center}
\newpage

\subsection{Compromiso de Notificación y Protección}
La Política invita a todo el personal a notificar peligros, errores y eventos, sin temor a represalias. Los reportes son confidenciales y no se usan para sancionar, salvo las excepciones de cultura justa descritas más abajo y en el Anexo 5. La Política también reconoce que cumplir el RAC 100 y el RAC 219 es obligatorio y que el SMS y sus procedimientos son de cumplimiento obligatorio para todo el personal.

\subsection{Aptitud para el Trabajo: Alcohol, Drogas y Fatiga}
\begin{itemize}
    \item \textbf{Cero tolerancia} al consumo de alcohol y sustancias psicoactivas (RAC 120): nadie opera, observa ni hace mantenimiento bajo sus efectos ni durante las horas previas. Si hay duda, no se trabaja y se informa al Gerente SMS.
    \item \textbf{Medicamentos:} quien tome un medicamento que cause sueño o afecte la atención lo informa al Jefe de Pilotos y no ocupa funciones críticas ese día.
    \item \textbf{Fatiga:} descanso de al menos 8 horas entre jornadas, jornada de vuelo y de mantenimiento fijada por el Jefe de Pilotos y el Gerente SMS, y derecho de cualquier persona a negarse a trabajar si está agotada. Quien lo haga no recibe ninguna sanción.
    \item \textbf{Funciones críticas} (piloto al mando, observador y técnico que libera aeronaves): antes de cada jornada, se declara con el FRAT (Anexo 6) que se está en condiciones de trabajar.
\end{itemize}

\subsection{Comunicación y Revisión de la Política}
\begin{enumerate}
    \item \textbf{Publicación:} el Gerente SMS entrega la Política firmada a todo el personal, y la deja visible en la oficina y en el repositorio digital. A clientes, proveedores y aseguradora se les entrega cuando lo soliciten o al contratar.
    \item \textbf{Comprensión:} en la inducción y en cada reentrenamiento se explica la Política y el personal firma que la recibió y la entendió (F-015). Se verifica con preguntas en el examen del Anexo 12.
    \item \textbf{Participación visible de la dirección:} el Gerente General la presenta personalmente en la reunión anual y firma las actas del SRB.
    \item \textbf{Eficacia:} en el SRB se revisa si el personal conoce la Política (resultados de evaluaciones y reportes) y si el medio de comunicación sirvió; si no, se cambia.
    \item \textbf{Revisión:} el Gerente General revisa la Política una vez al año y cuando cambie el alcance, y la vuelve a firmar. Se guardan las evidencias cinco años.
\end{enumerate}

\section{Objetivos de Seguridad Operacional}
Los objetivos del Sistema de Gestión de Seguridad Operacional (SMS) son fundamentales para guiar a la organización y establecer metas claras que permitan medir la eficacia y el desempeño del sistema. Para SKYFILMS, estos propósitos representan el compromiso de la alta dirección con la mitigación proactiva de los riesgos en todas nuestras operaciones con aeronaves no tripuladas (UAS).\\

\textbf{Objetivo General:}\\
Establecer, mantener y mejorar continuamente un Sistema de Gestión de Seguridad Operacional robusto que permita la identificación temprana de peligros, la mitigación efectiva de riesgos y la garantía de operaciones seguras, eficientes y en estricto cumplimiento normativo, manteniendo siempre los riesgos en un nivel aceptable.\\

\textbf{Objetivos Específicos:}\\
Para materializar nuestro objetivo general y asegurar la mejora continua del sistema, la organización ha definido los siguientes objetivos específicos:

\begin{itemize}
    \item \textbf{Mejora Continua:} Garantizar la implementación y efectividad del SMS a través de auditorías internas y revisiones periódicas por la dirección.

    \item \textbf{Cultura de Seguridad:} Fortalecer la cultura de seguridad operacional en todos los niveles mediante capacitación continua, talleres especializados y la promoción activa del reporte voluntario sin carácter punitivo.

    \item \textbf{Cumplimiento Normativo:} Asegurar el cumplimiento irrestricto de los estándares regulatorios aplicables, en particular las disposiciones del RAC 100, RAC 219 y las circulares informativas vigentes de la autoridad aeronáutica.

    \item \textbf{Gestión de Recursos Adecuados:} Proveer y gestionar de manera eficiente y oportuna los recursos humanos, técnicos y financieros necesarios para el correcto funcionamiento del SMS.
\end{itemize}



Con el fin de facilitar el seguimiento, control y la medición del desempeño de estas metas, SKYFILMS ha vinculado estos objetivos a Indicadores de Rendimiento de Seguridad Operacional (SPI) específicos y medibles. El desglose detallado de las metas, periodicidad, fórmulas de medición y planes de acción asociados a cada objetivo se encuentra documentado en el \textbf{Anexo 1 - Objetivos de Seguridad Operacional KPI's}.


\subsection{Objetivos, SPI y Metas}
Cada objetivo específico se relaciona con uno o más indicadores y metas, que están en el Anexo 1 y el Anexo 9. Para cumplir con el Programa Estatal de Seguridad Operacional, los objetivos se alinean con los del Estado (prevenir accidentes, reportar los eventos obligatorios y cumplir el RAC 100 y el RAC 219). El cumplimiento de los SPI hace parte de los objetivos específicos de Skyfilms. Los SPI cubren las operaciones certificadas: VLOS, vuelo nocturno, zona urbana y la flota Mavic 4 Pro y Pavo20 Pro II.

\subsection{Comunicación y Revisión de los Objetivos}
El Gerente General aprueba los objetivos. El Gerente SMS los comunica con la Política (mismo procedimiento) y los presenta al SRB una vez al año con los resultados de los SPI. Si un SPI no se cumple, se corrige o el objetivo se revisa y se deja constancia en el acta del SRB.

 \section{Cultura Justa en el Sistema de Gestión}
Para SKYFILMS, la Cultura Justa constituye la piedra angular sobre la cual se cimienta el Sistema de Gestión de Seguridad Operacional (SMS). De acuerdo con los principios del Documento 9859 de la OACI y los requisitos del RAC 219, la organización define la Cultura Justa como un entorno operativo en el cual se alienta, protege y capacita al personal (pilotos a distancia, observadores de espacio aéreo, personal de mantenimiento y directivos) para notificar de manera transparente cualquier información esencial relacionada con la seguridad operacional, garantizando que no serán objeto de castigos, represalias o acciones punitivas por sus acciones, omisiones o decisiones, siempre que estas correspondan a su experiencia y entrenamiento.\\

Este enfoque reconoce explícitamente que el error humano es una variable inevitable dentro de los sistemas tecnológicos y aeronáuticos complejos, incluidas las misiones comerciales de filmación y operaciones especiales con drones (UAS). Por consiguiente, la alta dirección de SKYFILMS se compromete formalmente a mantener un equilibrio ético y operativo entre dos ejes fundamentales:

\begin{itemize}
    \item \textbf{El Enfoque de Aprendizaje Organizacional:} Centrado en investigar y comprender qué falló en el sistema (defensas ausentes, fallos de equipos, factores ambientales o ambigüedad en los procedimientos) en lugar de buscar quién cometió el error, permitiendo extraer lecciones técnicas orientadas a la mejora continua.

    \item \textbf{La Exigencia de Responsabilidad (Accountability):} El establecimiento de fronteras claras e inequívocas que diferencien las fallas honestas del sistema de aquellos comportamientos que constituyen faltas disciplinarias o conductas inaceptables que comprometan la seguridad de las personas o bienes en superficie.
\end{itemize}

\subsection{Procedimiento de Cultura Justa}
 Ante la ocurrencia de cualquier evento de seguridad operacional, recepción de un reporte voluntario (VOR) u obligatorio (MOR), o hallazgo detectado en auditoría, el Gerente SMS y el Comité de Seguridad Operacional aplicarán formalmente el conducto regular analítico estructurado en las siguientes fases consecutivas:

 \definecolor{headerblue}{RGB}{180, 198, 231}
 
\renewcommand{\arraystretch}{1.4}
\setlength{\tabcolsep}{4pt}
 
\begin{table}[H]
\begin{adjustbox}{center}
\centering
\footnotesize
\begin{tabular}{|>{\centering\arraybackslash}m{1.1cm}|>{\centering\arraybackslash}m{2.9cm}|>{\RaggedRight\arraybackslash}m{5cm}|>{\centering\arraybackslash}m{3cm}|>{\centering\arraybackslash}m{2.3cm}|>{\centering\arraybackslash}m{3cm}|}
\hline
\rowcolor{gray!20}
\textbf{ÍTEM} & \textbf{ACTIVIDAD} & \textbf{DESCRIPCIÓN} & \textbf{RESPONSABLE} & \textbf{REGISTRO} & \textbf{FRECUENCIA} \\
\hline
 
1 &
Recepción del reporte &
Canalizar reportes sobre incidentes o situaciones inseguras, garantizando confidencialidad y libre expresión. &
Gerente SMS &
Formato de reporte &
Cada vez que se genere un reporte \\
\hline
 
2 &
Análisis preliminar &
Realizar un análisis inicial para determinar la necesidad de una investigación formal, según la naturaleza del evento. &
Gerente SMS &
Memo interno &
Según el caso \\
\hline
 
3 &
Investigación técnica &
Ejecutar una investigación objetiva que identifique causas y factores contribuyentes, enfocándose en el aprendizaje y no en la sanción. &
Gerente SMS &
Informe técnico &
Cada evento reportado \\
\hline
 
4 &
Clasificación de conducta &
Determinar la naturaleza del comportamiento (error, negligencia, violación) con base en criterios establecidos de cultura justa. &
Gerente SMS &
Matriz de clasificación &
Trimestral \\
\hline
 
5 &
Comité de revisión &
Evaluar los hallazgos y definir las acciones correspondientes: educativas, correctivas o disciplinarias si aplica. &
SAG -- Safety Action Group &
Acta de comité &
Según requerimiento \\
\hline
 
6 &
Comunicación de resultados &
Informar al personal implicado sobre las conclusiones del análisis, respetando la confidencialidad y promoviendo el aprendizaje. &
Gerente SMS &
Informe de retroalimentación &
Cuando se requiera \\
\hline
 
7 &
Seguimiento y verificación &
Supervisar el cumplimiento de las acciones definidas, verificando su impacto y cierre efectivo. &
Gerente SMS &
Consolidado de acciones &
Cuando aplique \\
\hline
 
8 &
Formación y sensibilización &
Implementar actividades de capacitación en cultura justa y seguridad operacional, dirigidas a todo el personal involucrado. &
Jefe de pilotos y gerente de SMS &
Registro de capacitaciones &
Semestral o según necesidad \\
\hline
 
9 &
Retroalimentación al sistema &
Documentar los aprendizajes obtenidos y proponer mejoras al SMS, MO o demás manuales aplicables. &
SAG &
Informe de mejora continua &
Semestral \\
\hline
 
\end{tabular}
\end{adjustbox}
\end{table}



\subsection{Protección de la Información y Tiempos de Gestión}
La protección de la identidad del reportante, la desidentificación de la información, las excepciones para divulgarla y los plazos de acuse, análisis, cierre y retroalimentación para cada tipo de reporte están definidos en el Anexo 5. La retroalimentación al personal se hace de forma individual al reportante y de forma general en el boletín de seguridad del trimestre y en la reunión de seguridad.

\textbf{Conductas que no se protegen (excepciones):} violación deliberada de normas, negligencia grave, falsificación de registros y operar bajo alcohol o sustancias psicoactivas. En esos casos, la decisión sobre la acción disciplinaria la toma el Gerente General con el informe del Gerente SMS.

\section{Rendición de Cuentas y Responsabilidades}
En SKYFILMS, la seguridad operacional es una responsabilidad compartida por todos los miembros de la organización. Sin embargo, existen roles y responsabilidades específicas asignadas a cada cargo para garantizar la implementación efectiva del Sistema de Gestión de Seguridad Operacional (SMS) y el cumplimiento de la normativa vigente.\\

\subsection{Gerente General:} El Gerente General es la máxima autoridad en SKYFILMS en lo que respecta a la seguridad operacional. Tiene la responsabilidad final en la conducción de los asuntos de certificación o aprobación de la organización, de las operaciones autorizadas y de todos los asuntos de seguridad operacional, incluyendo el establecimiento y la promoción de las Políticas del SMS. Además, es la autoridad final sobre las cuestiones de recursos humanos, técnicos y financieros necesarios para el rendimiento eficaz y eficiente del sistema de gestión.\\

\subsection{Gerente SMS:} El Gerente SMS es asignado para ejecutar y documentar todas las tareas de implementación, ejecución y mantenimiento del SMS en nombre del Gerente General. Sus responsabilidades incluyen:

\begin{itemize}
    \item Asegurarse de que los procesos del SMS se establezcan, implementen, ejecuten y mantengan.
    \item Promover la conciencia y una cultura positiva de seguridad operacional.
    \item Servir de enlace con la Aerocivil en cuestiones relacionadas con la seguridad operacional.
    \item Intercambiar lecciones valiosas aprendidas con otros operadores de UAS.
    \item Gestionar las investigaciones internas de incidentes y accidentes.
    \item Asegurarse de que se estén identificando los peligros y gestionando los riesgos de la operación del explotador UAS.
    \item Mantener la documentación de seguridad operacional.
    \item Organizar la instrucción en seguridad operacional.
\end{itemize}

\subsubsection{Perfil de Competencias}
La persona asignada como Gerente SMS demostrar conocimientos o formación en áreas del sector aeronáutico y contar con experiencia operacional aeronáutica de, al menos, doce (12) meses.\\

\textbf{NOTA:} De acuerdo con el RAC 1 “Definiciones”, “actividad aeronáutica” significa: “Conjunto de tareas y operaciones, directa o indirectamente relacionadas con el empleo de aeronaves civiles.”

\subsubsection{Designación, Dedicación y Recursos}
\begin{itemize}
    \item \textbf{Nombramiento:} el Gerente General designa por escrito al Gerente SMS, con una carta que indica su nombre, cargo y funciones: [nombre y cédula del Gerente SMS por diligenciar]. La carta se archiva como registro del SMS.
    \item \textbf{Línea jerárquica:} reporta directamente al Gerente General.
    \item \textbf{Dedicación y recursos:} dispone del tiempo necesario para sus funciones (mínimo medio día por semana, y más en las épocas de auditoría o evento), acceso al repositorio digital, apoyo de los demás cargos y presupuesto para capacitación.
    \item \textbf{Competencia:} formación en SMS (curso inicial y el curso de personal clave del Anexo 12), conocimiento de las operaciones UAS y experiencia aeronáutica de al menos doce meses, como se dijo arriba.
    \item \textbf{Conflicto de intereses:} no es Gerente General ni Jefe de Pilotos y no aprueba sus propias acciones. Si el evento le involucra directamente, el Gerente General designa otra persona para investigarlo.
\end{itemize}

\subsubsection{Autoridad para Detener Operaciones}
El Gerente General, el Gerente SMS y el Jefe de Pilotos pueden detener o cancelar un vuelo, una misión o una operación cuando haya un riesgo no aceptable. Cualquier piloto u observador puede detener su propio vuelo. Nadie recibe sanción por detener una operación por razones de seguridad. La reanudación la autoriza el Jefe de Pilotos cuando el riesgo vuelva a ser aceptable y, si la causa es una zona roja, el Gerente General.

\subsection{Jefe de Pilotos:} El Jefe de Pilotos es el encargado de supervisar y asegurar que todos los procedimientos establecidos en el manual de operaciones se cumplan según se haya establecido. Sus responsabilidades son:

\begin{itemize}
    \item Gestionar las solicitudes de autorización de vuelos UAS ante la Aerocivil.
    \item Garantizar que las operaciones del operador UAS se efectúen de manera segura, de acuerdo con el RAC 100 y demás normas vigentes aplicables.
    \item Asegurar la adecuada asignación y utilización de cada UA y de cada piloto UAS, de acuerdo con las competencias requeridas, para el desarrollo de la operación autorizada al operador UAS.
    \item Mantener un registro documental de las competencias de cada piloto UAS designado.
    \item Monitorizar la competencia de cada piloto UAS designado.
    \item Proponer las directrices y procedimientos estándar para las operaciones autorizadas a SKYFILMS, según corresponda, socializando, monitorizando y supervisando su cumplimiento por parte de cada piloto UAS designado y demás personal involucrado en las operaciones.
\end{itemize}

\subsection{Pilotos y Observadores:} Los pilotos y observadores son responsables de la operación directa de los UAS. Sus responsabilidades incluyen:

\begin{itemize}
    \item Realizar vuelos de manera segura y cumpliendo con la normativa.
    \item Mantener una comunicación efectiva con el control de tráfico aéreo y otros miembros del equipo.
    \item Reportar cualquier anomalía o condición insegura.
\end{itemize}

La asignación clara de responsabilidades y la rendición de cuentas son fundamentales para el éxito del SMS de SKYFILMS. Cada miembro del equipo es responsable de cumplir con sus obligaciones y de contribuir a la seguridad de las operaciones.

\section{Estructura Organizacional y Personal Clave}
SKYFILMS cuenta con una estructura organizacional jerárquica y funcionalmente definida para garantizar la eficiencia y seguridad de sus operaciones. Esta estructura facilita la comunicación, la toma de decisiones y la rendición de cuentas en todos los niveles. Además, dos grupos clave, el Grupo de Acción de Seguridad Operacional (SAG) y la Junta de Revisión de Seguridad Operacional (SRB), complementan esta estructura para una gestión integral de la seguridad.\\

En la cúspide se encuentra el \textbf{Gerente General}, quien tiene la autoridad final en todas las cuestiones de seguridad y es responsable del SMS. El \textbf{Gerente SMS} reporta directamente a él y se encarga de la implementación y mantenimiento del SMS.\\

El \textbf{Jefe de Pilotos} supervisa las operaciones de vuelo, asegurando el cumplimiento de procedimientos y normativas. Los \textbf{Pilotos} son responsables de la operación segura de los UAS y reportan al Jefe de Pilotos.\\

Para la gestión de riesgos y la mejora continua, SKYFILMS cuenta con dos grupos especializados:

\begin{itemize}
    \item La \textbf{Junta de Revisión de Seguridad Operacional (SRB)}, liderada por el Gerente General y conformada por el Gerente SMS y el Jefe de Pilotos, se enfoca en la gestión estratégica de la seguridad. Sus funciones incluyen la revisión del desempeño del SMS, la evaluación de riesgos de alto nivel, la asignación de recursos y la definición de directrices estratégicas para el SAG. Esta junta se reunirá con una periodicidad \textbf{anual} o cada que se identifique un peligro nuevo con un nivel de riesgo en la región \textbf{NO Tolerable} o \textbf{Tolerable} y cada reunión deberá tener un acta de realización emitida el \textbf{Gerente General}.
    \item El \textbf{Grupo de Acción de Seguridad Operacional (SAG)}, liderado por el Gerente SMS y conformado por los responsables de las áreas funcionales que influyen en la operación final, se encarga de la gestión táctica de la seguridad. Esto incluye la identificación de peligros, la evaluación y mitigación de riesgos, y la supervisión de la eficacia de las medidas de seguridad. Este grupo se reunirá cada que la Junta de Revisión de Seguridad Operacional lo solicite.
\end{itemize}

Esta estructura organizacional, que incluye los grupos SAG y SRB, asegura una clara definición de roles, facilita la comunicación y coordinación, y permite una gestión integral de la seguridad operacional en SKYFILMS. La tabla a continuación ilustra la jerarquía y las relaciones entre los diferentes cargos y grupos.

\begin{table}[H]
\centering
\renewcommand{\arraystretch}{1.35}
\begin{tabularx}{\textwidth}{|>{\RaggedRight\arraybackslash}p{4.4cm}|>{\RaggedRight\arraybackslash}X|>{\RaggedRight\arraybackslash}p{4.4cm}|}
\hline
\rowcolor[gray]{0.9}\textbf{Instancia} & \textbf{Función} & \textbf{Se relaciona con} \\ \hline
\rowcolor[HTML]{F4D6D0}Gerente General (Ejecutivo Responsable) & Máxima autoridad; toma las decisiones de seguridad operacional. & Preside el SRB \\ \hline
\rowcolor[HTML]{FFF2A8}Junta de Revisión de Seguridad Operacional (SRB) & Revisa el desempeño del SMS y decide sobre recursos y riesgos altos. & Gerente General; Gerente SMS; Jefe de Pilotos; SAG \\ \hline
\rowcolor[HTML]{FFF2A8}Gerente SMS (coordinador) & Administra el SMS y coordina las acciones. & SRB y SAG \\ \hline
Jefe de Pilotos & Supervisa la operación y aplica las acciones en vuelo. & SRB y SAG \\ \hline
\rowcolor[HTML]{FFF2A8}Grupo de Acción de Seguridad Operacional (SAG) & Ejecuta y da seguimiento a las acciones. & SRB; Gerente SMS; Jefe de Pilotos \\ \hline
\end{tabularx}
\caption{Estructura del SMS de Skyfilms.}
\end{table}


\subsection{Conflicto de Intereses}
En SKYFILMS como medida para evitar el conflicto de intereses se define que, el \textbf{Gerente SMS} no podrá ejercer en ningún momento el cargo de \textbf{Gerente General} ni \textbf{Jefe de pilotos}.


\subsection{Grupos de Seguridad: SRB y SAG}
\begin{longtable}{|>{\RaggedRight\arraybackslash}p{2.2cm}|>{\RaggedRight\arraybackslash}p{3.2cm}|>{\RaggedRight\arraybackslash}p{3.2cm}|>{\RaggedRight\arraybackslash}p{3.0cm}|>{\RaggedRight\arraybackslash}p{3.4cm}|}
\hline
\rowcolor[gray]{0.9}\textbf{Grupo} & \textbf{Integrantes} & \textbf{Frecuencia} & \textbf{Quórum} & \textbf{Qué hace y cómo se informa} \\ \hline
\endhead
SRB & Gerente General (preside), Gerente SMS, Jefe de Pilotos y Técnico UAS & Una vez cada seis meses y de inmediato ante un accidente, incidente grave o riesgo rojo & Gerente General y Gerente SMS & Revisa SPI, riesgos altos, recursos y auditorías. Se deja acta firmada y el Gerente SMS comunica las decisiones a todo el personal en cinco días hábiles \\ \hline
SAG & Gerente SMS (lidera), Jefe de Pilotos, Técnico UAS y un piloto & Cada trimestre y cuando haya un evento & Gerente SMS y dos integrantes más & Analiza reportes, riesgos y acciones; da seguimiento y escala al SRB lo que exceda su autoridad \\ \hline
\end{longtable}

\subsection{Niveles de Autoridad para Aceptar Riesgos}
\begin{longtable}{|>{\RaggedRight\arraybackslash}p{3cm}|>{\RaggedRight\arraybackslash}p{4.3cm}|>{\RaggedRight\arraybackslash}p{4.4cm}|>{\RaggedRight\arraybackslash}p{3.7cm}|}
\hline
\rowcolor[gray]{0.9}\textbf{Región} & \textbf{Quién decide} & \textbf{Qué se hace y en qué plazo} & \textbf{Revisión} \\ \hline
\endhead
Roja (no tolerable) & Gerente General & La operación se detiene de inmediato. Plan de acción en 24 horas. No se reanuda hasta que el riesgo residual sea amarillo o verde & Seguimiento semanal hasta el cierre \\ \hline
Amarilla con índice 4C, 4D, 4E, 5D, 5E o 3B & Gerente SMS, con informe al Gerente General en el SRB & Plan de acción en 5 días hábiles; controles puestos en 15 días hábiles & Cada mes \\ \hline
Amarilla (otros índices) & Gerente SMS, con visto bueno del Jefe de Pilotos & Plan de acción en 10 días hábiles; controles en 30 días hábiles & Cada trimestre \\ \hline
Verde (aceptable) & Gerente SMS & Se acepta y se vigila & Cada año \\ \hline
\end{longtable}

\section{Coordinación de Respuesta Ante Emergencias}
SKYFILMS, reconoce la importancia de contar con un plan de respuesta a emergencias (PRE) sólido y efectivo. Este PRE, en línea con las regulaciones colombianas, tiene como objetivo establecer los procedimientos y acciones necesarias para responder de manera oportuna y eficiente ante cualquier eventualidad que pueda afectar la seguridad de las operaciones, el personal, los equipos o el medio ambiente.

\subsection{Alcance}
Este PRE abarca todas las operaciones aéreas realizadas por SKYFILMS con sus UAS, así como las actividades desarrolladas en sus instalaciones y el bienestar de su personal. Esta sección debe ser conocida y consultada por todo el personal en caso de emergencia.

\subsection{Objetivos}

\begin{itemize}
    \item Proteger la vida e integridad de las personas involucradas en las operaciones de SKYFILMS.
    \item Minimizar los daños a los equipos, las instalaciones y el medio ambiente.
    \item Garantizar la continuidad de las operaciones en la medida de lo posible.
    \item Facilitar la comunicación y coordinación con las autoridades competentes y los servicios de emergencia.
    \item Establecer los procedimientos para la investigación de incidentes y la implementación de medidas correctivas.
\end{itemize}

\subsection{Marco Legal}
Este PRE se fundamenta en las siguientes normativas:

\begin{itemize}
    \item Reglamentos Aeronáuticos de Colombia (RAC) 100
    \item Reglamentos Aeronáuticos de Colombia (RAC) 219
\end{itemize}

\subsection{Estructura Organizacional de Respuesta a Emergencias}
La estructura de respuesta a emergencias de SKYFILMS se basa en los siguientes roles y responsabilidades:

\begin{itemize}
    \item \textbf{Gerente General:}
        \begin{itemize}
            \item Autoridad máxima en la toma de decisiones durante una emergencia.
            \item Responsable de la activación y ejecución del PRE.
            \item Enlace con las autoridades y los medios de comunicación.
        \end{itemize}
    \item \textbf{Gerente SMS:}
        \begin{itemize}
            \item Responsable de la coordinación general de la respuesta a la emergencia.
            \item Encargado de la comunicación interna y externa.
            \item Responsable de la investigación de incidentes y la implementación de medidas correctivas.
        \end{itemize}
    \item \textbf{Jefe de Pilotos:}
        \begin{itemize}
            \item Responsable de la seguridad de las operaciones de vuelo.
            \item Encargado de la coordinación de los pilotos y observadores durante una emergencia.
        \end{itemize}
    \item \textbf{Pilotos y Observadores:}
        \begin{itemize}
            \item Responsables de seguir los procedimientos de emergencia establecidos.
            \item Encargados de la comunicación con el Jefe de Pilotos y el Gerente SMS.
        \end{itemize}
\end{itemize}

\subsection{Procedimientos de Respuesta a Emergencias}

\subsubsection{Detección y Alerta}
\begin{itemize}
    \item Cualquier persona que presencie o tenga conocimiento de una emergencia deberá comunicarlo de inmediato al \textbf{Jefe de Pilotos} o, en su ausencia, al \textbf{Gerente SMS.}
    \item La comunicación se realizará utilizando los medios disponibles más rápidos (teléfono, radio, correo electrónico, etc.).
    \item Se proporcionará la siguiente información:
        \begin{itemize}
            \item Naturaleza de la emergencia
            \item Ubicación exacta
            \item Número de personas afectadas (si aplica)
            \item Gravedad de la situación
            \item Información de contacto del reportante
        \end{itemize}
\end{itemize}

\subsubsection{Evaluación Inicial}
\begin{itemize}
    \item El \textbf{Jefe de Pilotos} o el \textbf{Gerente SMS} evaluarán la información recibida y determinarán la magnitud de la emergencia.
    \item Se considerarán factores como la gravedad de la situación, el potencial de daños y la necesidad de activar el PRE.
\end{itemize}

\subsubsection{Activación del PRE}
\begin{itemize}
    \item Si la evaluación inicial indica que la situación lo amerita, el \textbf{Jefe de Pilotos} o el \textbf{Gerente SMS} informará al \textbf{Gerente General} y declarará la activación del PRE.
    \item Se informará de inmediato a todo el personal relevante y se iniciarán los procedimientos específicos correspondientes a la emergencia en cuestión.
\end{itemize}

\subsubsection{Notificación a Autoridades}
\begin{itemize}
    \item En caso de ser necesario, se notificará a las autoridades competentes, como la Aerocivil, la policía, los bomberos, los servicios de emergencia médica, etc.
    
    \item La notificación se realizará utilizando los números de contacto designados en el directorio telefónico del PRE \textbf{(Anexo 2 - Contactos de emergencia)}.
    \item Se proporcionará información detallada sobre la emergencia, incluyendo la ubicación, la naturaleza, la gravedad y el número de personas afectadas.
\end{itemize}

\subsubsection{Procedimiento de Comunicación Directa con Aerocivil}

Durante la ocurrencia de una emergencia operacional, incidente grave o accidente que involucre las aeronaves no tripuladas (UAS) de SKYFILMS, es imperativo establecer un flujo de comunicación oficial, ordenado y oportuno con la Aerocivil (Aerocivil). Este procedimiento garantiza la correcta articulación institucional y previene la divulgación de información inexacta o no confirmada.\\

\textbf{Personal Autorizado para la Comunicación}

Para mantener una línea de mando unificada y proteger la integridad de la información durante la crisis, se restringe la vocería oficial y el contacto directo con la autoridad aeronáutica exclusivamente a los siguientes cargos:

\begin{itemize}
    \item \textbf{Gerente General:} Como Ejecutivo Responsable ante la autoridad, encargado de las notificaciones corporativas y legales.

    \item \textbf{Gerente SMS:} Como enlace técnico principal para el reporte de la emergencia y el flujo de información de seguridad.

    \item \textbf{Jefe de Pilotos:} Como apoyo técnico para detallar aspectos operacionales específicos del evento ante los inspectores de la autoridad.
\end{itemize}

Ningún otro miembro del personal operativo, técnico o administrativo de SKYFILMS está autorizado para emitir declaraciones, entregar reportes preliminares o establecer contacto oficial con la Aerocivil durante una emergencia, debiendo escalar cualquier requerimiento a los cargos mencionados.\\

\textbf{Fines Específicos de la Comunicación}\\

El contacto de emergencia con la Aerocivil se activará de forma inmediata y prioritaria, ciñéndose estrictamente a los siguientes propósitos operativos y regulatorios:

\begin{itemize}
    \item \textbf{Coordinación Operativa:} Notificación y articulación técnica con el Grupo de Drones (UAS) de la Aerocivil respecto a la condición de la aeronave, el espacio aéreo afectado y las medidas iniciales de contención adoptadas por la empresa.

    \item \textbf{Búsqueda y Rescate (SAR):} Solicitud de apoyo, soporte y coordinación con los servicios de Búsqueda y Salvamento en caso de extravío de la aeronave en zonas críticas o afectación a personas en superficie.

    \item \textbf{Investigación de Accidentes:} Notificación inicial para la activación de los protocolos de preservación de evidencia y articulación con el grupo de Investigación de Accidentes e Incidentes de Aviación de la Aerocivil.
\end{itemize}


\textbf{Canales de Comunicación}
Las comunicaciones se realizarán única y exclusivamente a través de los canales de contacto de emergencia oficiales, directorios y líneas telefónicas actualizadas y publicadas en la página web oficial de la Aerocivil. Para garantizar una respuesta ágil, el Gerente SMS será responsable de verificar semestralmente la vigencia de estos números y mantenerlos actualizados en el \textbf{Anexo 2 - Contactos de Emergencia} del presente manual.\\

Para llevar a cabo la comunicación con Aerocivil se deberán ejecutar los siguientes pasos: 



\begin{table}[H]
\begin{adjustbox}{center}
\centering
\footnotesize
\begin{tabular}{|>{\centering\arraybackslash}m{0.06\linewidth}|m{0.55\linewidth}|>{\centering\arraybackslash}m{0.25\linewidth}|}
\hline
\rowcolor{gray!20}
\textbf{PASO} & \textbf{ACTIVIDAD} & \textbf{RESPONSABLE} \\
\hline

1 &
Identificar la situación de emergencia. &
Piloto UAS, Observador UAS \\
\hline

2 &
Reconocer la necesidad de contactar a la Aerocivil debido a emergencias, accidentes o incidentes. &
Jefe de Pilotos, Gerente SMS \\
\hline

3 &
Evaluar el canal de comunicación mas eficiente, considerando la urgencia para la toma de decisiones, el cual puede ser por medio telefónico o a través de correo electrónico. &
Jefe de Pilotos, Gerente SMS \\
\hline

4 &
Establecimiento de contacto inicial, Informando la naturaleza del mensaje y proporcionando detalles relevantes, como la ubicación y el número de identificación del UA. &
Gerente \\
\hline

5 &
Confirmación y registro de la comunicación, incluyendo la hora, la naturaleza del mensaje y cualquier acción tomada como resultado de la comunicación. &
Jefe de Pilotos, Gerente SMS \\
\hline

\end{tabular}
\end{adjustbox}
\caption{Procedimiento de Comunicación con Aerocivil}
\label{tab:responsabilidades_documentales}
\end{table}

\subsubsection{Respuesta Inicial}
\begin{itemize}
    \item Se implementarán las acciones inmediatas necesarias para controlar la situación, minimizar los daños y proteger a las personas y los bienes. Esto puede incluir:
    \begin{itemize}
        \item Suspensión de operaciones de vuelo.
        \item Aislamiento del área afectada.
        \item Evacuación de instalaciones.
        \item Prestación de primeros auxilios.
        \item Activación de equipos de respuesta a emergencias (si aplica).
    \end{itemize}
\end{itemize}

\subsubsection{Investigación}
\begin{itemize}
    \item Se llevará a cabo una investigación exhaustiva liderada por el Gerente SMS para determinar las causas de la emergencia y se documentarán los hallazgos.
    \item Se tomarán medidas correctivas y preventivas para evitar futuras ocurrencias y se informara por medio de boletines las lecciones aprendidas y/o retroalimentación.
\end{itemize}

\subsubsection{Comunicación con Medios}
\begin{itemize}
    \item En caso de que la emergencia tenga repercusión mediática, los únicos autorizados para comunicarse con los medios de comunicación son el \textbf{Gerente General, Gerente SMS y Jefe de Pilotos}.
    \item Se proporcionará información precisa, oportuna y transparente, evitando especulaciones o declaraciones que puedan perjudicar la imagen de SKYFILMS.
\end{itemize}

%\subsubsection{Apoyo a Colaborador y Familias}
%\begin{itemize}
  %  \item Se brindará apoyo a los colaboradores y sus familias en caso de que sean afectados por la emergencia.
    %\item Esto puede incluir asistencia psicológica, legal o económica, según sea necesario.
%\end{itemize}

\subsubsection{Revisión y Mejora}
\begin{itemize}
    \item El PRE será revisado y actualizado periódicamente, y siempre después de una emergencia, para asegurar su eficacia y adaptabilidad.
    \item Se realizarán simulacros y ejercicios de emergencia de manera regular para capacitar al personal y poner a prueba la efectividad del plan.
\end{itemize}

\subsection{Procedimiento Específico para Emergencias con Baterías (UAS)}
Considerando que las aeronaves no tripuladas operadas por SKYFILMS utilizan baterías recargables de alta densidad energética (principalmente baterías de litio), se establece el siguiente procedimiento específico para la prevención, control y respuesta ante emergencias asociadas a baterías, tales como sobrecalentamiento, incendio, explosión, fuga térmica o daño estructural.

\subsubsection{Identificación del Riesgo}
Se consideran situaciones de emergencia asociadas a baterías, entre otras:

\begin{itemize}
    \item Inflado, deformación o daño físico visible en la batería.

    \item Sobrecalentamiento anormal durante carga, descarga o posterior al vuelo.

    \item Emisión de humo, olor químico fuerte o presencia de chispa.

    \item Incendio o combustión espontánea.

    \item Impacto severo durante la operación.

    \item Caída del UAS con posible afectación del compartimiento de baterías.
\end{itemize}

\subsubsection{Procedimiento de Actuación en Caso de Emergencia}
En caso de identificarse una situación anormal con la batería:

\begin{enumerate}
    \item \textbf{Durante la operación en vuelo:}
        \begin{enumerate}
            \item Activar procedimiento de aterrizaje inmediato en zona segura.

            \item Priorizar área libre de personas, vehículos e infraestructura crítica.

            \item En caso de pérdida de control asociada a falla energética, activar procedimiento de contingencia establecido.
        \end{enumerate}
    
    \item \textbf{En tierra (durante carga o manipulación):}
        \begin{enumerate}
            \item Suspender inmediatamente el proceso de carga.

            \item Aislar la batería en un área despejada y ventilada.

            \item No manipular directamente la batería si presenta incendio activo.

            \item Utilizar extintor tipo ABC o agente adecuado para fuegos clase B/C.

            \item Notificar inmediatamente al responsable de seguridad operacional.

            \item Si el evento escala, activar el Plan de Respuesta a Emergencias y coordinar con los servicios de emergencia externos (Bomberos, autoridad local).
        \end{enumerate}
\end{enumerate}

\subsubsection{Disposición Final de Baterías Dañadas}
Las baterías que hayan presentado:

\begin{itemize}
    \item Incendio.

    \item Inflado severo.

    \item Daño estructural significativo.

    \item Falla eléctrica recurrente.
\end{itemize}
Deberán ser retiradas permanentemente del servicio y dispuestas conforme a la normativa ambiental vigente para residuos peligrosos.\\

\textbf{NOTA:} Este procedimiento forma parte integral del Plan de Respuesta Ante Emergencias de SKYFILMS y es de cumplimiento obligatorio para todo el personal operativo y técnico involucrado en la operación de UAS.

\subsection{Procedimientos Específicos: Pérdida de Control y Pérdida de Enlace}
Los pasos detallados están en el Manual de Operaciones; aquí se resumen:
\begin{enumerate}
    \item El piloto avisa en voz alta al observador y al equipo y mantiene la calma.
    \item Intenta recuperar el enlace o el control; si hay retorno automático, lo deja actuar y no vuela hacia personas.
    \item El observador sigue la aeronave con la vista y anota el último punto y la hora.
    \item Si cae o desaparece, se asegura el lugar, se pone a salvo a las personas y se llama al 123 si hay riesgo.
    \item El Jefe de Pilotos informa al Gerente SMS y este decide el reporte a la Aerocivil (Anexo 4).
\end{enumerate}

\subsection{Pruebas y Simulacros del Plan de Respuesta}
\begin{itemize}
    \item \textbf{Frecuencia:} un simulacro cada seis meses, que se alterna entre escenarios de batería en llamas, pérdida de control y lesiones a una persona.
    \item \textbf{Cómo se hace:} el Gerente SMS elige el escenario y lo avisa al personal con anticipación; el Jefe de Pilotos coordina; el piloto y el observador actúan según el plan.
    \item \textbf{Evaluación:} el Gerente SMS mide tiempos, comunicación y uso del equipo, y escribe un informe de una página con los resultados.
    \item \textbf{Mejora:} las debilidades pasan al Plan de Mejora Continua (Anexo 11) y se corrige el plan, si hace falta, en 30 días.
    \item \textbf{Distribución y coordinación:} el plan se entrega a todo el personal con el Manual. Cuando operemos con un cliente o una entidad externa (por ejemplo, un lugar con seguridad propia), el Gerente SMS acuerda antes con ella quién llama a quién.
    \item \textbf{Otras regiones:} para cada ciudad o región nueva, antes de operar el Gerente SMS agrega los contactos al Anexo 2 y confirma el hospital y los bomberos más cercanos.
\end{itemize}

\subsection{Información de Contactos}
SKYFILMS mantendrá actualizada la información de los contactos de emergencias para las ciudades principales de Colombia en el \textbf{Anexo 2 - Contactos de Emergencia} y para las operaciones fuera de las ciudades principales de Colombia, previo a la operación, el \textbf{Gerente SMS} sera responsable de contactar las autoridades del sitio.


\section{Documentación SMS}
El acceso a la documentación del Sistema de Gestión de Seguridad Operacional (SMS) de SKYFILMS es fundamental para garantizar la trazabilidad, la transparencia y la eficacia de nuestros procesos de seguridad. Toda la documentación que aporta al sistema de seguridad operacional lleva una trazabilidad y se almacena de manera digital, asegurando que la información esté actualizada y sea accesible.

\subsection{Tipos de Documentos}
Cada documento del SMS de SKYFILMS se identifica con un código único según su naturaleza, el cual incluye la versión del documento. Estos documentos se almacenan en una carpeta principal designada para el SMS, que contiene subcarpetas para cada manual principal: Manual SMS, Manual de Operaciones y Manual de Mantenimiento. Dentro de cada subcarpeta, se almacena la trazabilidad de los documentos específicos que cada manual requiere para su cumplimiento.

\subsubsection{Categorización y Nomenclatura Documental}
Para garantizar la protección de la información sobre seguridad operacional y fuentes conexas, todos los documentos del SMS de SKYFILMS se clasifican en tres niveles de acceso, identificados mediante una estructura alfanumérica estándar:\\

\begin{table}[H]
\begin{adjustbox}{center}
\begin{tabular}{|>{\centering\arraybackslash}m{0.3\linewidth}|>{\centering\arraybackslash}m{0.45\linewidth}|}
\hline
\rowcolor{gray!20}
\textbf{Nomenclatura} & \textbf{Categoría} \\
\hline

PUB & Publico \\
\hline

INT & Interno \\
\hline

CON & Confidencial \\
\hline

\end{tabular}
\end{adjustbox}
\caption{Nomenclatura y Categorización }
\label{tab:responsabilidades_documentales}
\end{table}

\begin{table}[H]
\begin{adjustbox}{center}
\begin{tabular}{|>{\centering\arraybackslash}m{0.3\linewidth}|>{\centering\arraybackslash}m{0.45\linewidth}|}
\hline
\rowcolor{gray!20}
\textbf{Rol} & \textbf{Nivel de Acceso} \\
\hline

Gerente General & CON \\
\hline

Gerente SMS & CON \\
\hline

Jefe de Pilotos & CON \\
\hline

Pilotos UAS & INT \\
\hline
Técnico UAS & INT \\
\hline

Clientes / Aseguradora / Proveedores & PUB \\
\hline

\end{tabular}
\end{adjustbox}
\caption{Nivel de Acceso por Rol}
\label{tab:responsabilidades_documentales}
\end{table}

Los niveles de acceso (PUB, INT y CON) dicen quién puede \textit{ver} cada documento. Quién lo elabora, revisa, aprueba o custodia (E, R, A, C, S) está en la matriz del Anexo 3. Son dos matrices distintas y no se mezclan.

\subsection{Control de la Documentación}
La gestión de la documentación y los registros operacionales en SKYFILMS garantiza que el personal tenga acceso a la información actualizada, que las operaciones se realicen bajo los estándares vigentes y que exista trazabilidad para la gestión de riesgos y auditorías de la Aerocivil, para ello se crea el \textbf{Anexo 3 - Procedimiento de Control de Documentos y Registros}.

\subsubsection{Directrices Generales de Gestión Documental}
Todos los documentos de SKYFILMS deben regirse por las siguientes directrices obligatorias:

\begin{itemize}
    \item \textbf{Identificación Estándar:} Todo documento debe incluir en su encabezado: Código de identificación, título del documento, logotipo de la empresa, fecha de emisión, número de versión y paginación (Ej. Página X de Y).
    
    \item \textbf{Aprobación y Autoría:} Cada documento o revisión debe tener claramente definido el responsable de elaboración (Autor), revisión y aprobación (Ej. elaborado por Gerente SMS, aprobado por Ejecutivo Responsable).
    
    \item \textbf{Control de Cambios:} Deberá existir un cuadro de control de revisiones y enmiendas (Listado de Páginas Efectivas), registrando la fecha de modificación y la justificación del cambio.
    
    \item \textbf{Trazabilidad y Retención:} Los registros operacionales (como bitácoras de vuelo y registros de mantenimiento) deben conservarse disponibles durante un periodo mínimo de 2 años en estado activo, y 3 años adicionales en archivo inactivo (archivo muerto).

    \item \textbf{Preservación de Evidencia:} En caso de incidente o accidente, el registro de los logs de vuelo, audios y videos de la estación de control deben preservarse y bloquearse de inmediato para proteger la información ante la Dirección Técnica de Investigación de Accidentes (DIACC).
\end{itemize}

\subsubsection{Lista Maestra de Documentos SMS}
El SDCPS (Sistema de recopilación y procesamiento de datos sobre seguridad operacional) de SKYFILMS se compone de los siguientes manuales, procedimientos y registros:

\begin{table}[H]
\begin{adjustbox}{center}
\begin{tabular}{|>{\centering\arraybackslash}m{0.08\linewidth}|m{0.70\linewidth}|>{\centering\arraybackslash}m{0.1\linewidth}|}
\hline
\rowcolor{gray!20}
\textbf{No} & \textbf{Documento} & \textbf{Nivel de Acceso} \\
\hline

1 & Manual del Sistema de Gestión de Seguridad Operacional & INT \\
\hline

2 & Manual de Operaciones UAS & INT  \\
\hline

3 & Plan de Respuesta ante Emergencias (PRE / ERP). & INT  \\
\hline

4 & Procedimientos Operacionales Normales y Anormales de aviación no tripulada. & INT \\
\hline

5 & Matriz de Identificación de Peligros y Evaluación de Riesgos & CON \\
\hline

6 & Actas de Comités de Seguridad Operacional (SRB estratégico y SAG táctico). & CON  \\
\hline

7 & Actas de reuniones de pilotos y socialización de boletines de seguridad. & INT \\
\hline

8 & Formularios de Reporte Obligatorio (MOR) y Reporte Voluntario (VOR). & INT \\
\hline

9 & Autorizaciones de vuelo UAS (emitidas por la Aerocivil) y planes de vuelo internos. & INT \\
\hline

10 & Libros de Vuelo por cada UA (Logbooks individuales por dron). & INT \\
\hline

11 & Bitácoras de vuelo por cada Piloto UAS. & INT \\
\hline

12 & Manual de Control de Mantenimiento (MCM) y formatos de registro de mantenimiento técnico. & INT \\
\hline

13 & Manuales técnicos del fabricante del UAS y de los Equipos Tecnológicos Asociados (ETA). & INT \\
\hline

14 & Registros de Instrucción, Capacitación y Evaluación de Competencias del personal en SMS. & INT \\
\hline

15 & Indicadores de Rendimiento de Seguridad Operacional (SPIs, KPIs y Metas SPT). & CON \\
\hline

16 & Informes de Auditorías Internas y Externas del SMS. & CON \\
\hline

\end{tabular}
\end{adjustbox}
\caption{Lista Maestra de Documentos SMS}
\label{tab:responsabilidades_documentales}
\end{table}

\subsubsection{Asignación de Responsabilidades Documentales}
Para evitar cruces de funciones (conflicto de intereses) exigidos por el RAC 100 y RAC 219, la responsabilidad de la administración, diligenciamiento y custodia de los documentos se divide de la siguiente manera:

\begin{table}[H]
\begin{adjustbox}{center}
\centering
\renewcommand{\arraystretch}{1.5}
\begin{tabular}{|m{0.3\textwidth}|m{0.7\textwidth}|}
\hline
\rowcolor{gray!20}
\textbf{Rol Responsable} & \textbf{Documentos y Registros a su cargo} \\ \hline
\textbf{Gerente SMS} & 
\textbullet~Manual del SMS (MSMS).\newline
\textbullet~Matriz de Peligros y Riesgos (SRM).\newline
\textbullet~Actas de reuniones de comités SMS (SRB / SAG).\newline
\textbullet~Formularios y bases de datos de Reportes VOR / MOR.\newline
\textbullet~Indicadores de Rendimiento (SPIs y KPIs) y análisis de tendencias.\newline
\textbullet~Registros de capacitación e instrucción SMS del personal.\newline
\textbullet~Informes de investigación de eventos y auditorías internas. \\ \hline
\textbf{Jefe de Pilotos UAS} & 
\textbullet~Manual de Operaciones (MO).\newline
\textbullet~Libros de vuelo individuales por cada aeronave no tripulada (UA).\newline
\textbullet~Bitácoras de control de horas de vuelo de los Pilotos UAS.\newline
\textbullet~Procedimientos operacionales, Instructivos y Listas de Chequeo (Checklists).\newline
\textbullet~Planificaciones y Autorizaciones de Vuelo UAS ante Aerocivil.\newline
\textbullet~Certificados de idoneidad, médicas y adiciones de su equipo de pilotos. \\ \hline
\end{tabular}
\end{adjustbox}
\caption{Asignación de Responsabilidades Documentales del SMS}
\label{tab:responsabilidades_documentales}
\end{table}




